\documentclass[12pt,a4paper]{article}
\usepackage{fontspec}
\setmainfont{Times New Roman}[Path=C:/Windows/Fonts/,UprightFont=times.ttf,BoldFont=timesbd.ttf,ItalicFont=timesi.ttf,BoldItalicFont=timesbi.ttf]
\usepackage{amsmath,amssymb,amsfonts}
\usepackage[left=1.5cm,right=1.5cm,top=1.5cm,bottom=1.7cm,headheight=17pt]{geometry}
\usepackage{enumitem,array,fancyhdr,booktabs}

\pagestyle{fancy}
\fancyhf{}
\lhead{\small\textbf{ĐÁP ÁN ĐỀ CƯƠNG CHƯƠNG IV -- NGUYÊN HÀM VÀ TÍCH PHÂN}}
\rhead{\small\textit{Biên soạn: Dương Gia Bảo}}
\rfoot{\small Trang \thepage}
\renewcommand{\headrulewidth}{0.4pt}
\fancypagestyle{firstpage}{\fancyhead{}\renewcommand{\headrulewidth}{0pt}}

\setlength{\parindent}{0pt}
\setlength{\parskip}{4pt}
\setlength{\emergencystretch}{1em}

\setlist[enumerate,1]{label=\textbf{Câu \arabic*:},leftmargin=*,itemsep=7pt,topsep=4pt}

\begin{document}
\thispagestyle{firstpage}

\noindent
\begin{minipage}[t]{0.36\textwidth}
    \centering
    {\fontsize{12}{14}\selectfont\textbf{ĐỀ CƯƠNG ÔN TẬP}}\\[4pt]
    \textit{Biên soạn: Dương Gia Bảo}\\[6pt]
    \setlength{\fboxrule}{0.8pt}
    \setlength{\fboxsep}{5pt}
    \fbox{\begin{minipage}{0.88\linewidth}
        \centering
        {\fontsize{12}{14}\selectfont\textbf{ĐÁP ÁN CHƯƠNG IV}}
    \end{minipage}}
\end{minipage}\hfill
\begin{minipage}[t]{0.62\textwidth}
    \centering
    {\fontsize{13}{15}\selectfont\textbf{ĐÁP ÁN \& HƯỚNG DẪN GIẢI CHI TIẾT}}\\[3pt]
    {\fontsize{12}{14}\selectfont\textbf{CHƯƠNG IV: NGUYÊN HÀM VÀ TÍCH PHÂN}}\\[3pt]
    {\fontsize{11}{13}\selectfont\textit{(GDPT 2018 -- Đầy đủ 50 câu trọng tâm)}}
\end{minipage}

\par\vspace{4pt}
\hrule height 0.6pt
\vspace{8pt}

\begin{center}
    {\fontsize{13}{15}\selectfont\textbf{BẢNG ĐÁP ÁN NHANH 50 CÂU}}
\end{center}

\begin{center}
\renewcommand{\arraystretch}{1.22}
\begin{tabular}{|c|c||c|c||c|c||c|c||c|c|}
\hline
\textbf{Câu} & \textbf{Đáp án} & \textbf{Câu} & \textbf{Đáp án} & \textbf{Câu} & \textbf{Đáp án} & \textbf{Câu} & \textbf{Đáp án} & \textbf{Câu} & \textbf{Đáp án} \\ \hline
\textbf{1} & \textbf{B} & \textbf{11} & \textbf{B} & \textbf{21} & $24$ & \textbf{31} & $41$ & \textbf{41} & $8$ \\ \hline
\textbf{2} & \textbf{D} & \textbf{12} & \textbf{C} & \textbf{22} & $114$ & \textbf{32} & $25$ & \textbf{42} & $4{,}5$ \\ \hline
\textbf{3} & \textbf{A} & \textbf{13} & \textbf{A} & \textbf{23} & $7$ & \textbf{33} & $21{,}4$ & \textbf{43} & $28$ \\ \hline
\textbf{4} & \textbf{C} & \textbf{14} & \textbf{B} & \textbf{24} & $11$ & \textbf{34} & $7$ & \textbf{44} & $12$ \\ \hline
\textbf{5} & \textbf{A} & \textbf{15} & \textbf{D} & \textbf{25} & $7$ & \textbf{35} & $11$ & \textbf{45} & $36$ \\ \hline
\textbf{6} & \textbf{D} & \textbf{16} & \textbf{C} & \textbf{26} & $12$ & \textbf{36} & $34$ & \textbf{46} & $36$ \\ \hline
\textbf{7} & \textbf{B} & \textbf{17} & \textbf{B} & \textbf{27} & $22$ & \textbf{37} & $2$ & \textbf{47} & $9$ \\ \hline
\textbf{8} & \textbf{C} & \textbf{18} & \textbf{A} & \textbf{28} & $192$ & \textbf{38} & $4$ & \textbf{48} & $6$ \\ \hline
\textbf{9} & \textbf{A} & \textbf{19} & \textbf{D} & \textbf{29} & $300$ & \textbf{39} & $362{,}5$ & \textbf{49} & $21$ \\ \hline
\textbf{10} & \textbf{D} & \textbf{20} & \textbf{C} & \textbf{30} & $45$ & \textbf{40} & $126$ & \textbf{50} & $255$ \\ \hline
\end{tabular}
\end{center}

\vspace{8pt}
\hrule height 0.6pt
\vspace{8pt}

\begin{center}
    {\fontsize{13}{15}\selectfont\textbf{HƯỚNG DẪN GIẢI CHI TIẾT}}
\end{center}

\subsection*{PHẦN I. TRẮC NGHIỆM NHIỀU LỰA CHỌN (CÂU 1 -- 20)}

\begin{enumerate}
\item \textbf{Đáp án B.}  
Theo định nghĩa nguyên hàm, hàm số $F(x)$ là một nguyên hàm của $f(x)$ trên khoảng $K$ khi và chỉ khi $F'(x) = f(x)$ với mọi $x \in K$.

\item \textbf{Đáp án D.}  
$\int (4x^3 + 3x^2 - 2)\,\mathrm{d}x = x^4 + x^3 - 2x + C$.

\item \textbf{Đáp án A.}  
Trên $(0; +\infty)$, ta có $\sqrt{x} = x^{\frac{1}{2}}$, suy ra $\int \left(\dfrac{3}{x} + x^{\frac{1}{2}}\right)\,\mathrm{d}x = 3\ln x + \dfrac{2}{3}x^{\frac{3}{2}} + C = 3\ln x + \dfrac{2}{3}x\sqrt{x} + C$.

\item \textbf{Đáp án C.}  
Ta có $\int \left(\dfrac{1}{\sin^2 x} - 3\cos x\right)\,\mathrm{d}x = -\cot x - 3\sin x + C$.

\item \textbf{Đáp án A.}  
$\int (2^x + 3\mathrm{e}^x)\,\mathrm{d}x = \dfrac{2^x}{\ln 2} + 3\mathrm{e}^x + C$.

\item \textbf{Đáp án D.}  
$F(x) = \int (6x^2 - 4x + 1)\,\mathrm{d}x = 2x^3 - 2x^2 + x + C$.  
Vì $F(0) = 3 \implies C = 3 \implies F(2) = 2(8) - 2(4) + 2 + 3 = 13$.

\item \textbf{Đáp án B.}  
$v(t) = \int (4t + 2)\,\mathrm{d}t = 2t^2 + 2t + C$. Vì $v(0) = 0 \implies C = 0$.  
Tại $t = 3$, ta có $v(3) = 2(3)^2 + 2(3) = 24\text{ m/s}$.

\item \textbf{Đáp án C.}  
Theo công thức Newton--Leibniz: $\int_a^b f(x)\,\mathrm{d}x = F(b) - F(a) = -3 - 7 = -10$.

\item \textbf{Đáp án A.}  
$\int_{-1}^5 f(x)\,\mathrm{d}x = \int_{-1}^2 f(x)\,\mathrm{d}x + \int_2^5 f(x)\,\mathrm{d}x \implies \int_2^5 f(x)\,\mathrm{d}x = 12 - 5 = 7$.

\item \textbf{Đáp án D.}  
$I = 5\int_1^3 f(x)\,\mathrm{d}x + 3\int_1^3 g(x)\,\mathrm{d}x = 5(4) + 3(-1) = 17$.

\item \textbf{Đáp án B.}  
$I = \left.\left(x^4 + 3\mathrm{e}^x\right)\right|_0^1 = (1 + 3\mathrm{e}) - (0 + 3) = 3\mathrm{e} - 2$.

\item \textbf{Đáp án C.}  
$I = \left.(\sin 2x + x)\right|_0^{\frac{\pi}{4}} = \left(\sin\dfrac{\pi}{2} + \dfrac{\pi}{4}\right) - 0 = 1 + \dfrac{\pi}{4}$.

\item \textbf{Đáp án A.}  
$I = \int_{-1}^0 (-2x)\,\mathrm{d}x + \int_0^2 2x\,\mathrm{d}x = \left.(-x^2)\right|_{-1}^0 + \left.x^2\right|_0^2 = 1 + 4 = 5$.

\item \textbf{Đáp án B.}  
$I = \left.(2\ln|x|)\right|_1^2 = 2\ln 2 - 2\ln 1 = 2\ln 2 = \ln 4$.

\item \textbf{Đáp án D.}  
Công thức tính diện tích hình phẳng giới hạn bởi hai đồ thị $y = f(x)$, $y = g(x)$ và $x = a$, $x = b$ ($a < b$) là $S = \int_a^b |f(x) - g(x)|\,\mathrm{d}x$.

\item \textbf{Đáp án C.}  
Hoành độ giao điểm với $Ox$: $4 - x^2 = 0 \iff x = \pm 2$.  
$S = \int_{-2}^2 (4 - x^2)\,\mathrm{d}x = 2\left.\left(4x - \dfrac{x^3}{3}\right)\right|_0^2 = 2\left(8 - \dfrac{8}{3}\right) = \dfrac{32}{3}$.

\item \textbf{Đáp án B.}  
Hoành độ giao điểm: $x^2 - 1 = 1 - x^2 \iff 2x^2 = 2 \iff x = \pm 1$.  
$S = \int_{-1}^1 [(1 - x^2) - (x^2 - 1)]\,\mathrm{d}x = \int_{-1}^1 (2 - 2x^2)\,\mathrm{d}x = 2\left.\left(2x - \dfrac{2}{3}x^3\right)\right|_0^1 = \dfrac{8}{3}$.

\item \textbf{Đáp án A.}  
Diện tích thiết diện là $S(x) = \left(\sqrt{3x + 1}\right)^2 = 3x + 1$.  
Thể tích vật thể: $V = \int_0^2 (3x + 1)\,\mathrm{d}x = \left.\left(\dfrac{3}{2}x^2 + x\right)\right|_0^2 = 6 + 2 = 8$.

\item \textbf{Đáp án D.}  
$V = \pi\int_0^1 \left(\mathrm{e}^x\right)^2\,\mathrm{d}x = \pi\int_0^1 \mathrm{e}^{2x}\,\mathrm{d}x = \pi\left.\left(\dfrac{1}{2}\mathrm{e}^{2x}\right)\right|_0^1 = \dfrac{\pi(\mathrm{e}^2 - 1)}{2}$.

\item \textbf{Đáp án C.}  
Lượng nước bơm được từ $t = 1$ đến $t = 4$ là $\int_1^4 (3t^2 + 2t)\,\mathrm{d}t = \left.\left(t^3 + t^2\right)\right|_1^4 = (64 + 16) - (1 + 1) = 78\text{ lít}$.
\end{enumerate}

\subsection*{PHẦN II. TRẢ LỜI NGẮN (CÂU 21 -- 50)}

\begin{enumerate}[start=21]
\item \textbf{Đáp số: $24$.}  
$F(x) = x^3 - 4x^2 + 5x + C$. Vì $F(1) = 1 - 4 + 5 + C = 6 \implies C = 4$.  
Vậy $F(4) = 4^3 - 4\cdot 4^2 + 5\cdot 4 + 4 = 24$.

\item \textbf{Đáp số: $114$.}  
Trên $(0; +\infty)$, $F(x) = \int \left(2x + 6x^{-\frac{1}{2}}\right)\,\mathrm{d}x = x^2 + 12\sqrt{x} + C$.  
$F(1) = 1 + 12 + C = 10 \implies C = -3$. Suy ra $F(9) = 81 + 12(3) - 3 = 114$.

\item \textbf{Đáp số: $7$.}  
$f(x) = 2x - 1 + \dfrac{3}{x} \implies F(x) = x^2 - x + 3\ln x + C$ (vì $x > 0$).  
$F(1) = 1 - 1 + 0 + C = 4 \implies C = 4$.  
Do đó $F(\mathrm{e}) = \mathrm{e}^2 - \mathrm{e} + 3\ln\mathrm{e} + 4 = \mathrm{e}^2 - \mathrm{e} + 7 \implies c = 7$.

\item \textbf{Đáp số: $11$.}  
$F(x) = -4\cos x + 2\sin x + C$. Vì $F(0) = -4 + 0 + C = 3 \implies C = 7$.  
Khi đó $F(\pi) = -4\cos\pi + 2\sin\pi + 7 = 4 + 0 + 7 = 11$.

\item \textbf{Đáp số: $7$.}  
Ta có $f(x) = \tan^2 x = \dfrac{1}{\cos^2 x} - 1 \implies F(x) = \tan x - x + C$.  
$F(0) = 2 \implies C = 2$. Suy ra $F\!\left(\dfrac{\pi}{4}\right) = \tan\dfrac{\pi}{4} - \dfrac{\pi}{4} + 2 = 3 - \dfrac{\pi}{4}$.  
Vậy $a = 3, b = 4 \implies a + b = 7$.

\item \textbf{Đáp số: $12$.}  
$f(x) = 2\mathrm{e}^x + 1 \implies F(x) = 2\mathrm{e}^x + x + C$.  
$F(0) = 2 + C = 5 \implies C = 3$. Vậy $F(\ln 3) = 2\cdot 3 + \ln 3 + 3 = 9 + \ln 3$, suy ra $a = 9, b = 3 \implies a + b = 12$.

\item \textbf{Đáp số: $22$.}  
$f(x) = \int (3x^2 - 4x + 2)\,\mathrm{d}x = x^3 - 2x^2 + 2x + C$.  
Vì đồ thị qua $M(2; 11)$ nên $f(2) = 8 - 8 + 4 + C = 11 \implies C = 7$.  
Do đó $f(3) = 27 - 18 + 6 + 7 = 22$.

\item \textbf{Đáp số: $192$.}  
$h(t) = \int (1{,}5t + 20)\,\mathrm{d}t = 0{,}75t^2 + 20t + C$.  
Vì $h(0) = 45 \implies C = 45$. Sau $6$ năm: $h(6) = 0{,}75(36) + 20(6) + 45 = 27 + 120 + 45 = 192\text{ cm}$.

\item \textbf{Đáp số: $300$.}  
$C(x) = \int (0{,}06x^2 - 1{,}2x + 15)\,\mathrm{d}x = 0{,}02x^3 - 0{,}6x^2 + 15x + C_0$.  
Với $C(0) = 80 \implies C_0 = 80$. Khi $x = 20$:  
$C(20) = 0{,}02(8000) - 0{,}6(400) + 15(20) + 80 = 160 - 240 + 300 + 80 = 300$ triệu đồng.

\item \textbf{Đáp số: $45$.}  
Vận tốc: $v(t) = \int (-10)\,\mathrm{d}t = -10t + 20$. Vật đạt độ cao lớn nhất khi $v(t) = 0 \iff t = 2\text{ s}$.  
Độ cao: $h(t) = \int (-10t + 20)\,\mathrm{d}t = -5t^2 + 20t + 25$ (do $h(0) = 25$).  
Độ cao cực đại: $h(2) = -5(4) + 20(2) + 25 = 45\text{ m}$.

\item \textbf{Đáp số: $41$.}  
Đặt $A = \int_0^4 f(x)\,\mathrm{d}x$ và $B = \int_0^4 g(x)\,\mathrm{d}x$. Ta có hệ:
\[
\begin{cases} 2A + B = 19 \\ A - B = 2 \end{cases} \iff \begin{cases} A = 7 \\ B = 5 \end{cases}.
\]
Khi đó $I = 3A + 4B = 3(7) + 4(5) = 41$.

\item \textbf{Đáp số: $25$.}  
\begin{align*}
\int_0^2 \dfrac{2x + 5}{x + 1}\,\mathrm{d}x &= \int_0^2 \left(2 + \dfrac{3}{x + 1}\right)\,\mathrm{d}x = \left.\bigl[2x + 3\ln(x + 1)\bigr]\right|_0^2 = 4 + 3\ln 3.
\end{align*}
Suy ra $a = 4, b = 3 \implies P = 4^2 + 3^2 = 25$.

\item \textbf{Đáp số: $21{,}4$.}  
\begin{align*}
I &= \int_1^4 \left(x^{\frac{3}{2}} + 3\right)\,\mathrm{d}x = \left.\left(\dfrac{2}{5}x^{\frac{5}{2}} + 3x\right)\right|_1^4 \\
&= \left(\dfrac{2}{5}\cdot 32 + 12\right) - \left(\dfrac{2}{5} + 3\right) = \dfrac{62}{5} + 9 = 12{,}4 + 9 = 21{,}4.
\end{align*}

\item \textbf{Đáp số: $7$.}  
\begin{align*}
\int_0^{\frac{\pi}{2}} (2\sin x + 3\cos x - 1)\,\mathrm{d}x &= \left.(-2\cos x + 3\sin x - x)\right|_0^{\frac{\pi}{2}} \\
&= \left(0 + 3 - \dfrac{\pi}{2}\right) - (-2 + 0 - 0) = 5 - \dfrac{\pi}{2}.
\end{align*}
Vậy $a = 5, b = 2 \implies a + b = 7$.

\item \textbf{Đáp số: $11$.}  
Ta có $x^2 - 2x \le 0$ trên $[0; 2]$ và $x^2 - 2x \ge 0$ trên $[2; 3]$.  
\begin{align*}
I &= \int_0^2 (2x - x^2)\,\mathrm{d}x + \int_2^3 (x^2 - 2x)\,\mathrm{d}x \\
&= \left.\left(x^2 - \dfrac{x^3}{3}\right)\right|_0^2 + \left.\left(\dfrac{x^3}{3} - x^2\right)\right|_2^3 = \dfrac{4}{3} + \left(0 - \left(-\dfrac{4}{3}\right)\right) = \dfrac{8}{3}.
\end{align*}
Vậy $a = 8, b = 3 \implies a + b = 11$.

\item \textbf{Đáp số: $34$.}  
\begin{align*}
I &= \int_{-1}^2 (x^2 + 3)\,\mathrm{d}x + \int_2^4 (4x - 1)\,\mathrm{d}x \\
&= \left.\left(\dfrac{x^3}{3} + 3x\right)\right|_{-1}^2 + \left.\left(2x^2 - x\right)\right|_2^4 = \left(\dfrac{8}{3} + 6 + \dfrac{1}{3} + 3\right) + (28 - 6) = 12 + 22 = 34.
\end{align*}

\item \textbf{Đáp số: $2$.}  
$\int_0^m (3x^2 - 4x + 4)\,\mathrm{d}x = m^3 - 2m^2 + 4m = 8 \iff (m - 2)(m^2 + 4) = 0 \iff m = 2$.

\item \textbf{Đáp số: $4$.}  
Ta có $v(t) = t^2 - 4t + 3 = 0 \iff t = 1$ hoặc $t = 3$.  
Tổng quãng đường đi được trong $4$ giây đầu tiên là:
\begin{align*}
S &= \int_0^4 |t^2 - 4t + 3|\,\mathrm{d}t \\
&= \int_0^1 (t^2 - 4t + 3)\,\mathrm{d}t - \int_1^3 (t^2 - 4t + 3)\,\mathrm{d}t + \int_3^4 (t^2 - 4t + 3)\,\mathrm{d}t \\
&= \dfrac{4}{3} + \dfrac{4}{3} + \dfrac{4}{3} = 4\text{ m}.
\end{align*}

\item \textbf{Đáp số: $362{,}5$.}  
Quãng đường bằng tổng diện tích ba miền (tam giác, hình chữ nhật, hình thang vuông) dưới đường gấp khúc $OABC$:
\begin{align*}
S &= \int_0^{30} v(t)\,\mathrm{d}t \\
&= \dfrac{1}{2}\cdot 10\cdot 15 + 15\cdot (25 - 10) + \dfrac{15 + 10}{2}\cdot (30 - 25) \\
&= 75 + 225 + 62{,}5 = 362{,}5\text{ m}.
\end{align*}

\item \textbf{Đáp số: $126$.}  
$E = \int_0^6 (12 + 6t - 0{,}75t^2)\,\mathrm{d}t = \left.\left(12t + 3t^2 - 0{,}25t^3\right)\right|_0^6 = 72 + 108 - 54 = 126\text{ kWh}$.

\item \textbf{Đáp số: $8$.}  
Hoành độ giao điểm với $Ox$: $x^3 - 4x = 0 \iff x \in \{-2; 0; 2\}$.  
Vì $y = x^3 - 4x$ là hàm số lẻ nên:
\[
S = 2\int_0^2 (4x - x^3)\,\mathrm{d}x = 2\left.\left(2x^2 - \dfrac{x^4}{4}\right)\right|_0^2 = 2(8 - 4) = 8.
\]

\item \textbf{Đáp số: $4{,}5$.}  
Hoành độ giao điểm: $x^2 - 4x + 5 = -x + 5 \iff x^2 - 3x = 0 \iff x = 0$ hoặc $x = 3$.  
$S = \int_0^3 (3x - x^2)\,\mathrm{d}x = \left.\left(\dfrac{3}{2}x^2 - \dfrac{1}{3}x^3\right)\right|_0^3 = \dfrac{27}{2} - 9 = \dfrac{9}{2} = 4{,}5$.

\item \textbf{Đáp số: $28$.}  
Từ đồ thị, $S_1 = \int_{-2}^1 f(x)\,\mathrm{d}x = 11$ và $S_2 = -\int_1^4 f(x)\,\mathrm{d}x = 6 \implies \int_{-2}^4 f(x)\,\mathrm{d}x = 11 - 6 = 5$.  
Do đó:
\[
I = 2\int_{-2}^4 f(x)\,\mathrm{d}x + \int_{-2}^4 3\,\mathrm{d}x = 2(5) + 3\cdot (4 - (-2)) = 10 + 18 = 28.
\]

\item \textbf{Đáp số: $12$.}  
Hai parabol cắt nhau tại $x = 0$ và $x = 6$. Trên $[0; 6]$, $-\dfrac{1}{6}x^2 + 2x \ge \dfrac{1}{6}x^2$.  
\[
S = \int_0^6 \left(2x - \dfrac{1}{3}x^2\right)\,\mathrm{d}x = \left.\left(x^2 - \dfrac{1}{9}x^3\right)\right|_0^6 = 36 - 24 = 12\text{ dm}^2.
\]

\item \textbf{Đáp số: $36$.}  
Chọn hệ trục tọa độ với đỉnh parabol ở $(0; 9)$ và hai chân hầm ở $(-6; 0), (6; 0)$.  
Phương trình parabol là $y = 9 - \dfrac{1}{4}x^2$. Diện tích mặt cắt ngang cửa hầm:
\[
S = \int_{-6}^6 \left(9 - \dfrac{1}{4}x^2\right)\,\mathrm{d}x = 2\left.\left(9x - \dfrac{x^3}{12}\right)\right|_0^6 = 2(54 - 18) = 72\text{ m}^2.
\]
Chi phí làm lưới chắn là $72 \times 0{,}5 = 36$ triệu đồng.

\item \textbf{Đáp số: $36$.}  
Tam giác vuông cân có cạnh huyền $c = 2\sqrt{x^2 + 2x}$ thì diện tích thiết diện là $S(x) = \dfrac{1}{4}c^2 = x^2 + 2x$.  
Thể tích vật thể:
\[
V = \int_1^4 (x^2 + 2x)\,\mathrm{d}x = \left.\left(\dfrac{x^3}{3} + x^2\right)\right|_1^4 = \left(\dfrac{64}{3} + 16\right) - \left(\dfrac{1}{3} + 1\right) = 21 + 15 = 36.
\]

\item \textbf{Đáp số: $9$.}  
Diện tích thiết diện tam giác đều cạnh $6 - 2x$ là $S(x) = \dfrac{\sqrt{3}}{4}(6 - 2x)^2$.  
\[
V = \dfrac{\sqrt{3}}{4}\int_0^3 (6 - 2x)^2\,\mathrm{d}x = \dfrac{\sqrt{3}}{4}\left.\left(-\dfrac{(6 - 2x)^3}{6}\right)\right|_0^3 = \dfrac{\sqrt{3}}{4} \cdot \dfrac{216}{6} = 9\sqrt{3}\text{ m}^3 \implies a = 9.
\]

\item \textbf{Đáp số: $6$.}  
$V = \pi\int_0^2 \left(\sqrt{x^3 + 1}\right)^2\,\mathrm{d}x = \pi\int_0^2 (x^3 + 1)\,\mathrm{d}x = \pi\left.\left(\dfrac{x^4}{4} + x\right)\right|_0^2 = \pi(4 + 2) = 6\pi \implies \dfrac{V}{\pi} = 6$.

\item \textbf{Đáp số: $21$.}  
$V = \pi\int_0^3 (x + 1)^2\,\mathrm{d}x = \pi\left.\left(\dfrac{(x + 1)^3}{3}\right)\right|_0^3 = \pi\left(\dfrac{64}{3} - \dfrac{1}{3}\right) = 21\pi \implies \dfrac{V}{\pi} = 21$.

\item \textbf{Đáp số: $255$.}  
Dung tích của thùng rượu là thể tích khối tròn xoay:
\begin{align*}
V &= \pi\int_{-3}^3 \left(4 - \dfrac{1}{9}x^2\right)^2\,\mathrm{d}x = 2\pi\int_0^3 \left(16 - \dfrac{8}{9}x^2 + \dfrac{1}{81}x^4\right)\,\mathrm{d}x \\
&= 2\pi\left.\left(16x - \dfrac{8}{27}x^3 + \dfrac{1}{405}x^5\right)\right|_0^3 = 2\pi\left(48 - 8 + \dfrac{3}{5}\right) = \dfrac{406\pi}{5} \approx 255{,}1\text{ lít}.
\end{align*}
Làm tròn đến hàng đơn vị thu được $255$ lít.
\end{enumerate}

\end{document}
